\chapter{Geodesic rays and test curves}
\index{geodesic ray}\index{test curve}\index{Legendre transform!ray--test-curve}
\label{chap:rays-test-curves}

Test curves turn the convex geometry of geodesic rays into a calculus of
singularity types.  The transform is pointwise Legendre duality, but its
analytic content is much stronger: model envelopes are exactly what make
the transform produce weak geodesics rather than arbitrary
subgeodesics.  This chapter develops that bridge and then records the
operations needed for the inverse-limit construction of
\Cref{chap:transcendental-metrics}.

Throughout, \(X\) is a connected compact K\"ahler manifold,
\(\theta\) represents a big class, and
\(\phi\in\PSH(X,\theta)_{>0}\) is a model potential.  We write
\[
 \mathcal E(X,\theta;\phi)
 :=\{\varphi\in\PSH(X,\theta)\mid P_\theta[\varphi]=\phi\}
\tag{52.1}\label{eq:rays-relative-full-mass}
\]
for the relative full-mass class.  Additive normalization is not included
in this notation.  For the ray formulas below we choose the representative
of the model type with
\[
                              \sup_X\phi=0.
\]
This choice is indispensable in \eqref{eq:rays-linear-supremum}; replacing
\(\phi\) by \(\phi+C\) translates every potential by \(C\), whereas
translating the test-curve parameter by \(C\) adds \(Ct\) to the ray.  The
two operations have different types and will never be identified.

We use the squared-norm radial variable
\[
                              t=-\log|z|^2.
\]
Thus the Legendre term is \(+t\tau\).  If a source uses
\(-\log|z|\), then its radial parameter is \(t/2\); one must change the
parameter and the normalization of generic Lelong numbers together.  No
factor of two is inserted into the formulas of this chapter.

\section{Weak geodesics and rays}

Let \(S_{a,b}=\{z\in\mathbf C\mid a<\Re z<b\}\) and let
\(p_1:X\times S_{a,b}\to X\) be projection.

\begin{definition}[Subgeodesic]
\label{def:rays-subgeodesic}
A family \((u_t)_{a<t<b}\subset\PSH(X,\theta)\) is a subgeodesic if
\[
                 U(x,z):=u_{\Re z}(x)
\tag{52.2}\label{eq:rays-complexification}
\]
is \(p_1^*\theta\)-psh.  On a closed interval we additionally require
the prescribed endpoint values in \(L^1\).
\end{definition}

For each \(x\), the map \(t\mapsto u_t(x)\) is convex or identically
\(-\infty\).  Hence a subgeodesic has rigid one-sided limits.  Adding an
affine function of \(t\), or changing \(\theta\) by \(\ddc f\) and
replacing every \(u_t\) by \(u_t-f\), preserves the notion.

\begin{definition}[Weak geodesic segment]
\label{def:rays-geodesic-segment}
For \(u_0,u_1\in\PSH(X,\theta)\), their weak geodesic is the Perron
envelope of all subgeodesics with boundary values at most \(u_0,u_1\):
\[
 u_t=\uscsup\{v_t\mid (v_s)_{0<s<1}\text{ is such a subgeodesic}\}.
\tag{52.3}\label{eq:rays-perron-geodesic}
\]
\end{definition}

The envelope can be \(-\infty\) without a common lower barrier.  In the
relative full-mass class the barrier is supplied by the rooftop.

\begin{theorem}\label{thm:rays-relative-geodesic}
If \(u_0,u_1\in\mathcal E(X,\theta;\phi)\), then
\eqref{eq:rays-perron-geodesic} is a subgeodesic with the correct
\(L^1\)-boundary values, every \(u_t\) lies in
\(\mathcal E(X,\theta;\phi)\), and each restriction to a subinterval is
again the Perron geodesic between its endpoints.
\end{theorem}

\begin{proof}[Imported inputs and Perron deductions]
We use two deep results from the relative pluripotential theory developed in
the supplied ray chapter of \cite{XiaBook}: relative rooftop stability in
\(\mathcal E(X,\theta;\phi)\), and the Perron boundary theorem for endpoints
in that class.  We spell out how they enter.

After adding endpoint constants, assume \(u_0,u_1\leq\phi\).  Relative
rooftop stability gives
\[
                    u_0\wedge_\theta u_1\in
                    \mathcal E(X,\theta;\phi),
\]
so the constant family at the rooftop is a Perron competitor.  The
maximum principle for \(p_1^*\theta\)-psh functions proves that the usc
Perron envelope is a subgeodesic.  To identify the boundary at \(0\),
replace \(u_0\) by \(u_0\wedge_\theta(u_1+C)\), use the affine barrier,
and let \(C\to\infty\); relative full-mass rooftop convergence gives
\(L^1\)-convergence to \(u_0\).  The argument at \(1\) is symmetric.
Balayage on a smaller strip proves the final assertion.
The rooftop convergence and boundary passage are the imported analytic
steps; the maximum-principle and balayage deductions above are formal once
those inputs are available.
\end{proof}

\begin{definition}[Geodesic ray]
\label{def:rays-geodesic-ray}
A geodesic ray from \(\phi\) is a family
\(\ell=(\ell_t)_{t\geq0}\subset\mathcal E(X,\theta;\phi)\) with
\(\ell_0=\phi\), whose restriction to each compact interval is a weak
geodesic.  Denote the set of such rays by
\(\mathcal R(X,\theta;\phi)\).
\end{definition}

\begin{lemma}\label{lem:rays-linear-supremum}
Every ray has linear supremum:
\[
                    \sup_X\ell_t=t\sup_X\ell_1.
\tag{52.4}\label{eq:rays-linear-supremum}
\]
\end{lemma}

\begin{proof}
Put
\(a(t):=\sup_{\{\phi>-\infty\}}(\ell_t-\phi)\).  The maximum principle
for a relative weak geodesic segment says that \(a\) is affine on every
compact interval on which the ray is geodesic.  Since \(\ell_0=\phi\),
we have \(a(0)=0\), and compatibility of the affine functions on nested
intervals gives \(a(t)=ta(1)\).  We now use the model-potential supremum
identity from the supplied ray chapter of \cite{XiaBook},
\[
 \sup_{\{\phi>-\infty\}}(u-\phi)=\sup_Xu
\quad(u\preceq\phi).
\tag{52.5}\label{eq:rays-model-supremum}
\]
This identity uses the chosen normalization \(\sup_X\phi=0\).  Applying it to
\(u=\ell_t\) and \(u=\ell_1\) proves
\eqref{eq:rays-linear-supremum}.
\end{proof}

\section{Test curves}

\begin{definition}[Test curve]
\label{def:rays-test-curve}
A test curve \(\Gamma\) consists of
\(\Gamma_{\max}\in\mathbf R\) and a map
\[
 (-\infty,\Gamma_{\max})\longrightarrow\PSH(X,\theta),
 \qquad \tau\longmapsto\Gamma_\tau ,
\]
such that:
\begin{enumerate}
\item \(\tau\mapsto\Gamma_\tau\) is decreasing and concave;
\item every level is a fixed point of the \(P\)-envelope,
      \(P_\theta[\Gamma_\tau]=\Gamma_\tau\).  Positive mass is \emph{not}
      imposed here --- \Cref{def:complex-p-envelope} reserves the term
      ``model potential'' for the positive-mass case, and reading it into
      this clause would make \Cref{lem:rays-test-positive-levels} vacuous.
      Positivity is imposed only at \(\Gamma_{-\infty}\), by
      \eqref{eq:rays-test-minus-infinity}, and deduced at interior levels in
      that lemma;
\item
      \[
      \Gamma_{-\infty}
      :=\uscsup_{\tau<\Gamma_{\max}}\Gamma_\tau
      \quad\text{satisfies}\quad
      \int_X\theta_{\Gamma_{-\infty}}^n>0.
      \tag{52.6}\label{eq:rays-test-minus-infinity}
      \]
\end{enumerate}
Write \(\mathcal{TC}(X,\theta;\phi)\) when
\(\Gamma_{-\infty}=\phi\).
\end{definition}

Extend a test curve by
\[
 \Gamma_\tau=-\infty\quad(\tau>\Gamma_{\max}),\qquad
 \Gamma_{\Gamma_{\max}}
 =\inf_{\tau<\Gamma_{\max}}\Gamma_\tau .
\tag{52.7}\label{eq:rays-test-endpoint}
\]
This is an extended-real convention, not a claim that the endpoint has
positive mass.

\begin{lemma}\label{lem:rays-test-positive-levels}
If \(\Gamma\) is a test curve, then
\[
                   \int_X\theta_{\Gamma_\tau}^n>0
 \quad\text{for every }\tau<\Gamma_{\max}.
\tag{52.8}\label{eq:rays-test-positive-levels}
\]
\end{lemma}

\begin{proof}
Choose \(\tau_0<\tau\) so negative that the mass of
\(\Gamma_{\tau_0}\) is positive; this is possible by increasing
continuity at \(-\infty\).  Choose
\(\tau'\in(\tau,\Gamma_{\max})\) and \(s\in(0,1)\) with
\(\tau=s\tau_0+(1-s)\tau'\).  Concavity gives
\[
 \Gamma_\tau\geq s\Gamma_{\tau_0}+(1-s)\Gamma_{\tau'}.
\]
Multilinearity and mass monotonicity then give
\[
 \int_X\theta_{\Gamma_\tau}^n
 \geq s^n\int_X\theta_{\Gamma_{\tau_0}}^n>0.
\]
\end{proof}

A test curve is bounded if \(\Gamma_\tau=\phi\) for all sufficiently
negative \(\tau\).  It has finite energy if
\[
 \mathbf E^\phi(\Gamma)
 :=\Gamma_{\max}M_\theta(\phi)
 +\int_{-\infty}^{\Gamma_{\max}}
       \bigl(M_\theta(\Gamma_\tau)-M_\theta(\phi)\bigr)\,\rd \tau
 >-\infty .
\tag{52.9}\label{eq:rays-test-energy}
\]
The integral is meaningful because the integrand is nonpositive and the
level-mass function is continuous on the interior.

\section{Legendre transform and the RWN correspondence}

\begin{definition}[Legendre transforms]
\label{def:rays-legendre}
For a test curve \(\Gamma\), define, for \(t>0\),
\[
                 \Gamma_t^*
 :=\sup_{\tau\in\mathbf R}(t\tau+\Gamma_\tau).
\tag{52.10}\label{eq:rays-forward-legendre}
\]
There is no usc regularization in this formula.  Set
\(\Gamma_0^*:=\phi\) by convention; formula
\eqref{eq:rays-forward-legendre} need not hold at \(t=0\).
For a ray \(\ell\), define
\[
                 \ell_\tau^*
 :=\inf_{t>0}(\ell_t-t\tau).
\tag{52.11}\label{eq:rays-inverse-legendre}
\]
\end{definition}

The two directions of the transform are displayed in
\cref{fig:rays-test-curves-legendre}.  The profiles in the drawing are scalar
slices at one fixed point; the definitions above are function-valued.

\begin{figure}[htbp]
\centering
\begin{tikzpicture}[x=0.9cm,y=1cm,>=Stealth]
  % A scalar slice of a decreasing concave test curve.
  \draw[-{Stealth[length=2mm]},thin] (0,0)--(4.15,0)
    node[right,font=\small] {$\tau$};
  \draw[-{Stealth[length=2mm]},thin] (0,0)--(0,3.15);
  \node[rotate=90,font=\small] at (-0.45,1.7) {$\Gamma_\tau(x)$};
  \draw[very thick]
    plot[smooth] coordinates
      {(0.25,2.75) (1.15,2.62) (2.05,2.2) (2.85,1.42) (3.55,0.35)};
  \draw[dashed,thin] (3.55,0)--(3.55,2.95);
  \node[below,font=\scriptsize] at (3.55,0) {$\Gamma_{\max}$};
  \node[font=\scriptsize,align=center] at (1.85,3.38)
    {decreasing, concave profile};

  % A scalar slice of the corresponding convex ray.
  \draw[-{Stealth[length=2mm]},thin] (7,0)--(11.15,0)
    node[right,font=\small] {$t>0$};
  \draw[-{Stealth[length=2mm]},thin] (7,0)--(7,3.15);
  \node[rotate=90,font=\small] at (6.55,1.7) {$\ell_t(x)$};
  \draw[very thick]
    plot[smooth] coordinates
      {(7.25,0.35) (8.15,0.48) (9.05,0.9) (9.85,1.62) (10.55,2.72)};
  \node[font=\scriptsize,align=center] at (9.15,3.38)
    {convex ray profile};

  \draw[-{Stealth[length=2.2mm]},thick]
    (4.4,2.18)--(6.72,2.18)
    node[midway,above,font=\scriptsize,align=center]
      {$\ell_t=\sup_\tau(t\tau+\Gamma_\tau)$};
  \draw[-{Stealth[length=2.2mm]},thick]
    (6.72,0.92)--(4.4,0.92)
    node[midway,below,font=\scriptsize,align=center]
      {$\Gamma_\tau=\inf_{t>0}(\ell_t-t\tau)$};
\end{tikzpicture}
\caption{The Legendre correspondence at the level of a scalar slice.  The
upper arrow is \eqref{eq:rays-forward-legendre}; the lower arrow is
\eqref{eq:rays-inverse-legendre} together with biconjugacy.  The two vertical
axes have independent scales, so the drawing records convexity and the
direction of the transforms rather than numerical values.}
\label{fig:rays-test-curves-legendre}
\end{figure}

The extended-real arithmetic is classical:
\(-\infty+c=-\infty\) for finite \(c\), and an infimum equal to
\(-\infty\) is retained.  We never subtract two infinite values.

\begin{proposition}\label{prop:rays-legendre-well-defined}
For a test curve \(\Gamma\), each \(\Gamma_t^*\) is
\(\theta\)-psh and the family is a subgeodesic on \(t>0\).
For a ray \(\ell\), each \(\ell_\tau^*\) is
\(\theta\)-psh or identically \(-\infty\), and
\[
 \ell_\tau^*\in\PSH(X,\theta)
 \quad\Longleftrightarrow\quad
 \tau\leq\sup_X\ell_1.
\tag{52.12}\label{eq:rays-inverse-domain}
\]
\end{proposition}

\begin{proof}
For fixed \(x\), the functions in
\eqref{eq:rays-forward-legendre} are affine in \(t\), hence their
supremum is convex.  On the half-plane \(\Re z>0\), the family
\(t\tau+\Gamma_\tau\) is locally uniformly bounded above by
\(t\Gamma_{\max}\).  Kiselman's minimum principle and the psh supremum
lemma give joint upper semicontinuity, hence
\(\theta\)-psh slices.  For \eqref{eq:rays-inverse-legendre}, Kiselman's
principle applies to the \(S^1\)-invariant complexification of the ray.
The threshold follows from
\(\sup_X\ell_t=t\sup_X\ell_1\).
\end{proof}

\begin{theorem}\label{thm:rays-rwn}
Legendre transform gives inverse bijections
\[
 \mathcal{TC}(X,\theta;\phi)
 \ \xleftrightarrow{\ \ *\ \ }\
 \mathcal R(X,\theta;\phi).
\tag{52.13}\label{eq:rays-rwn}
\]
It restricts to bounded objects and to finite-energy objects.  For a
finite-energy test curve,
\[
                       \mathbf E^\phi(\Gamma)
                       =\mathbf E^\phi(\Gamma^*),
\tag{52.14}\label{eq:rays-rwn-energy}
\]
where the right side is the radial Monge--Amp\`ere energy slope.
\end{theorem}

\begin{proof}[Source boundary and normalization check]
The analytic assertions that the forward transform of a maximal test curve
is a weak geodesic, that the inverse transform of a weak geodesic has model
levels, and that their energies agree are the main theorems of
\cite[Theorems~1.1--1.2]{RWN14}.  Their relative full-mass formulation is
proved in the supplied test-curve chapter of \cite{XiaBook}.  These results
use comparison principles and Monge--Amp\`ere energy estimates that are not
reproved here.

We verify that the imported correspondence has exactly the conventions of
this chapter.  At every point outside the common pluripolar locus, the
closed concave function \(\tau\mapsto\Gamma_\tau(x)\) satisfies
Fenchel--Moreau biconjugacy with the sign \(+t\tau\):
\[
 \Gamma_\tau=\inf_{t>0}(\Gamma_t^*-t\tau).
\tag{52.15}\label{eq:rays-biconjugate}
\]
Canonical pluripolar extension makes this an identity on \(X\).  Conversely,
the imported inverse transform has effective domain
\(\tau\leq\sup_X\ell_1\) by
\eqref{eq:rays-inverse-domain}; its \(-\infty\) limit is
\(\ell_0=\phi\).  Thus both transforms land in the spaces displayed in
\eqref{eq:rays-rwn}, and biconjugacy makes them inverse.

For bounded curves, the source energy identity is obtained from finite
Legendre maxima and Riemann sums for
\eqref{eq:rays-test-energy}.  The source then truncates below \(-A\) and
uses monotone convergence.  With our squared-norm variable
\(t=-\log|z|^2\), its slope is precisely the right side of
\eqref{eq:rays-rwn-energy}; no factor of two occurs.  This proves the
normalization and the stated bounded and finite-energy restrictions.
\end{proof}

\begin{remark}[Sign convention]
\label{rem:rays-legendre-sign}
Our transform uses \(t\tau+\Gamma_\tau\).  Some sources parameterize the
annulus by \(t=-\log|z|\) and write the dual variable with the opposite
sign.  Formulae in this book are always interpreted through
\eqref{eq:rays-forward-legendre}--\eqref{eq:rays-inverse-legendre}.
\end{remark}

\section{Relative and \texorpdfstring{\(\mathcal I\)}{I}-model objects}

\begin{definition}[\(\mathcal I\)-model test curve]
\label{def:rays-i-model-test}
A test curve is \(\mathcal I\)-model if
\[
                 P_\theta[\Gamma_\tau]_{\mathcal I}
                 =\Gamma_\tau
 \quad\text{for every }\tau<\Gamma_{\max}.
\tag{52.16}\label{eq:rays-i-model-level}
\]
Write
\(\mathcal E^{\NA}(X,\theta;\phi)\) for these curves with
\(\Gamma_{-\infty}=\phi\), and
\[
 \PSH^{\NA}(X,\theta)_{>0}
 :=\bigcup_{\phi\ {\rm model}}
      \mathcal E^{\NA}(X,\theta;\phi).
\tag{52.17}\label{eq:rays-positive-na-test-curves}
\]
\end{definition}

The notation anticipates \Cref{chap:transcendental-metrics}.
Every level, including the \(-\infty\) level, is then
\(\mathcal I\)-model.  A geodesic ray is maximal precisely when its
inverse Legendre transform is \(\mathcal I\)-model.

For any test curve, the levelwise operation
\[
 (P_\theta[\Gamma]_{\mathcal I})_\tau
 :=P_\theta[\Gamma_\tau]_{\mathcal I}
\tag{52.18}\label{eq:rays-test-i-envelope}
\]
produces an \(\mathcal I\)-model test curve.  The two properties of the
envelope that this uses are monotonicity and concavity under convex
combinations,
\[
 P_\theta[(1-s)u+sv]_{\mathcal I}
 \geq(1-s)P_\theta[u]_{\mathcal I}
       +sP_\theta[v]_{\mathcal I}.
\tag{52.19}\label{eq:rays-i-concavity}
\]
The verification is recorded as a statement rather than left in passing,
because \Cref{chap:transcendental-metrics} invokes it at every K\"ahler
perturbation and because each test-curve axiom consumes a different
property of the envelope.

\begin{lemma}[Levelwise \(\mathcal I\)-envelope of a test curve]
\label{lem:rays-levelwise-i-envelope}
Let \(\Gamma\) be a test curve in the big class \(\theta\).  Then
\eqref{eq:rays-test-i-envelope} is an \(\mathcal I\)-model test curve with
the same terminal parameter \(\Gamma_{\max}\), and its \(-\infty\) level is
\[
 \bigl(P_\theta[\Gamma]_{\mathcal I}\bigr)_{-\infty}
 =P_\theta[\Gamma_{-\infty}]_{\mathcal I}.
\]
\end{lemma}

\begin{proof}
Write \(\Lambda_\tau:=P_\theta[\Gamma_\tau]_{\mathcal I}\).

Monotonicity of the \(\mathcal I\)-envelope carries the decrease of
\(\tau\mapsto\Gamma_\tau\) to the decrease of \(\tau\mapsto\Lambda_\tau\),
which is the first axiom of \Cref{def:rays-test-curve}.  For concavity, let
\(\tau=(1-s)\tau_1+s\tau_2\) with \(\tau_1,\tau_2<\Gamma_{\max}\) and
\(s\in(0,1)\).  Concavity of \(\Gamma\) gives
\(\Gamma_\tau\geq(1-s)\Gamma_{\tau_1}+s\Gamma_{\tau_2}\), so monotonicity
followed by \eqref{eq:rays-i-concavity} yields
\[
 \Lambda_\tau
 \geq P_\theta\bigl[(1-s)\Gamma_{\tau_1}+s\Gamma_{\tau_2}\bigr]_{\mathcal I}
 \geq(1-s)\Lambda_{\tau_1}+s\Lambda_{\tau_2}.
\]
Every level is \(\mathcal I\)-model by idempotence of the
\(\mathcal I\)-envelope, which is \eqref{eq:rays-i-model-level}.

For the positive-mass axiom \eqref{eq:rays-test-minus-infinity}, the
envelope dominates its argument, \(\Lambda_\tau\geq\Gamma_\tau\), so mass
monotonicity gives
\(M_\theta(\Lambda_\tau)\geq M_\theta(\Gamma_\tau)\), and the right side is
positive for every \(\tau<\Gamma_{\max}\) by
\Cref{lem:rays-test-positive-levels}.  No level below \(\Gamma_{\max}\) is
therefore lost.  No level above it is created either: there
\(\Gamma_\tau=-\infty\) by the extension convention
\eqref{eq:rays-test-endpoint}, and the envelope of the identically
\(-\infty\) function is again identically \(-\infty\).  Hence the terminal
parameter is unchanged.

The \(-\infty\) level is identified by increasing continuity of the
\(\mathcal I\)-envelope along the increasing family
\((\Gamma_\tau)\) as \(\tau\downarrow-\infty\).
\end{proof}

\begin{proposition}\label{prop:rays-kahler-transition}
If \(\omega\) is a smooth positive \((1,1)\)-form, then
\[
 P_{\theta+\omega}[\Gamma]_{\mathcal I}
 \in\PSH^{\NA}(X,\theta+\omega)_{>0}.
\tag{52.20}\label{eq:rays-kahler-transition}
\]
For \(\omega\geq\omega'\), the transitions compose:
\[
 P_{\theta+\omega}
 [P_{\theta+\omega'}[\Gamma]_{\mathcal I}]_{\mathcal I}
 =P_{\theta+\omega}[\Gamma]_{\mathcal I}.
\tag{52.21}\label{eq:rays-transition-cocycle}
\]
\end{proposition}

\begin{proof}
Apply the \(\mathcal I\)-envelope levelwise.  Its monotonicity and
concavity preserve the test-curve axioms, while adding the positive form
\(\omega\) preserves positive mass at every interior level by mass
monotonicity.  Thus \eqref{eq:rays-kahler-transition} is a positive-mass
\(\mathcal I\)-model test curve.  If \(\omega\geq\omega'\), both sides of
\eqref{eq:rays-transition-cocycle} are the normalized
\(\mathcal I\)-model representative in the class \(\theta+\omega\) of
the same multiplier-ideal type.  Idempotence of the
\(\mathcal I\)-envelope therefore gives their equality level by level.
\end{proof}

\section{Operations and monotone limits}

For test curves in possibly different big classes, define the pointwise
order by
\[
 \Gamma\leq\Gamma'
 \quad\Longleftrightarrow\quad
 \Gamma_{\max}\leq\Gamma'_{\max}
 \quad\text{and}\quad
 \Gamma_\tau\preceq_P\Gamma'_\tau
 \quad(\tau<\Gamma_{\max}),
\tag{52.22}\label{eq:rays-test-order}
\]
with endpoint behavior understood through the extension
\eqref{eq:rays-test-endpoint}.
For normalized model levels, \(P\)-comparison is ordinary pointwise order
of the model representatives.  When the levels are \(\mathcal I\)-model,
the constructions below use their \(\mathcal I\)-envelopes in addition.
This does not identify the two ambient equivalence relations: outside the
\(\mathcal I\)-good positive-mass locus one only has
\(\preceq_P\Rightarrow\preceq_{\mathcal I}\), as recorded in
\Cref{chap:complex-prerequisites}.

The sum is the modelized sup-convolution
\[
 (\Gamma+\Gamma')_\tau
 :=P_{\theta+\theta'}
 \left[\uscsup_{\delta\in\mathbf R}
       \bigl(\Gamma_\delta+\Gamma'_{\tau-\delta}\bigr)\right]_{\mathcal I}.
\tag{52.23}\label{eq:rays-test-sum}
\]
The obstacle in brackets is \(\mathcal I\)-good in the positive-mass
regime by the addition theorem for test curves in the supplied operations
chapter of \cite{XiaBook}.  Consequently its \(P\)- and
\(\mathcal I\)-envelopes agree.  Formula
\eqref{eq:rays-test-sum} records the \(\mathcal I\)-envelope explicitly
because that is the form compatible with K\"ahler transition.
Translation shifts the parameter:
\[
              (\Gamma+C)_\tau:=\Gamma_{\tau-C},
\qquad
              (\Gamma+C)_{\max}=\Gamma_{\max}+C.
\tag{52.24}\label{eq:rays-test-translation}
\]
The maximum is obtained by taking the least \(\mathcal I\)-model test
curve dominating both inputs.  These definitions are arranged so that
Legendre transform gives
\[
 \begin{split}
 (\Gamma+\Gamma')^*&=\Gamma^*+\Gamma'^*,\\
 (\Gamma+C)^*_t&=\Gamma_t^*+Ct,\\
 (\Gamma\vee\Gamma')^*_t&=\Gamma_t^*\vee\Gamma_t'^*.
 \end{split}
\tag{52.25}\label{eq:rays-operations-legendre}
\]

\begin{proposition}\label{prop:rays-test-limits}
The following operations remain in the class of \(\mathcal I\)-model
test curves.
\begin{enumerate}
\item A nonempty family with
      \(\sup_i\Gamma^i_{\max}<+\infty\) has a regularized supremum.
      A countable subfamily has the same supremum.
\item An increasing net with the same upper bound on
      \(\Gamma^i_{\max}\) converges to that regularized supremum.
\item A decreasing net with
      \(\inf_i\Gamma^i_{\max}>-\infty\) has a pointwise infimum test
      curve.
\item For \(\lambda>0\), the rescaled curve is
      \[
      (\lambda\Gamma)_\tau
          :=\lambda\Gamma_{\tau/\lambda}
      \quad\text{in the class }\lambda\theta.
      \tag{52.26}\label{eq:rays-test-rescaling}
      \]
\end{enumerate}
All four constructions commute with K\"ahler transition
\eqref{eq:rays-kahler-transition}.
\end{proposition}

\begin{proof}[Limit safeguards]
At each level, ordinary Choquet reduction gives a countable cofinal
subfamily for a uniformly bounded-above psh envelope.  Apply it first
at rational \(\tau\), then use concavity to recover all real levels.
The rational-level family is decreasing in \(\tau\); taking its closed
concave extension supplies the test curve and shows that no
parameter-dependent choice of a countable subfamily remains.

For a decreasing net, the lower bound on the terminal parameters keeps a
nontrivial common interval.  At any fixed interior level, the positive-mass
infimum identity \cite[Proposition~2.8]{Xia23Operations} identifies the
pointwise limit with its \(\mathcal I\)-model projection.  Concavity passes
to the infimum and the endpoint is supplied by the extended-curve
convention.  Rescaling is checked by substituting
\(\rho=\tau/\lambda\); the current scales as
\(\lambda(\theta+\ddc\Gamma_\rho)\), so positivity and the model condition
are preserved.

Compatibility with K\"ahler transition follows from the same
\(\mathcal I\)-envelope identities for finite maxima, admissible suprema,
and decreasing infima.  Without the displayed bounds on
\(\Gamma_{\max}\), the purported limit may be identically infinite or
have empty effective domain.
\end{proof}

\section{Filtrations and projective test configurations}

Assume now that \(L\) is big, \(h\) is a smooth Hermitian metric, and
\(\theta=c_1(L,h)\).  Let
\[
 R(X,L)=\bigoplus_{m\geq0}H^0(X,mL)
\]
and let \(\mathcal F\) be a decreasing, left-continuous,
multiplicative, linearly bounded filtration.  Thus, for some \(C>0\),
\[
 \mathcal F^{Cm}H^0(X,mL)=0,\qquad
 \mathcal F^{-Cm}H^0(X,mL)=H^0(X,mL).
\tag{52.27}\label{eq:rays-linear-filtration}
\]
For \(\tau\in\mathbf R\), set
\[
 \Gamma_{\tau}^{\mathcal F,m}
 :=\uscsup\left\{\log|s|_{h^m}^2\ \middle|\
 s\in\mathcal F^{m\tau}H^0(X,mL),\
 |s|_{h^m}^2\leq1\right\}.
\tag{52.28}\label{eq:rays-bergman-test}
\]
Multiplicativity and Fekete's lemma give the asymptotic obstacle
\[
 g_\tau^{\mathcal F}
 :=\uscsup_{m\geq1}\frac1m\Gamma_{\tau}^{\mathcal F,m}.
\tag{52.29}\label{eq:rays-filtration-obstacle}
\]
Define
\[
       \Gamma_\tau^{\mathcal F}
       :=P_\theta[g_\tau^{\mathcal F}].
\tag{52.30}\label{eq:rays-filtration-test-curve}
\]

\begin{theorem}\label{thm:rays-filtration-interface}
The family \(\Gamma^{\mathcal F}\) is a bounded
\(\mathcal I\)-model test curve, and its Legendre transform is a maximal
weak geodesic ray.  If \(L\) is ample and \(\mathcal F\) is induced by a
smooth relatively ample projective test configuration, this ray is the
Phong--Sturm ray.  No extension of this identification to an arbitrary
normal test configuration is asserted here.  For an ample \(L\), the outer
\(P_\theta\) in
\eqref{eq:rays-filtration-test-curve} is unnecessary by
\cite[Proposition~7.11]{RWN14}; no such removal is asserted here for an
arbitrary big \(L\).
\end{theorem}

\begin{proof}[Algebra-to-analysis map]
Choose sections, normalized by their sup norms, with
\[
 s\in\mathcal F^{m\tau}H^0(X,mL),
 \qquad
 s'\in\mathcal F^{m'\tau'}H^0(X,m'L).
\]
Multiplicativity puts \(ss'\) in filtration level
\(m\tau+m'\tau'\).  Taking logarithms and envelopes gives
\[
 \frac{m}{m+m'}\frac1m\Gamma_\tau^{\mathcal F,m}
 +\frac{m'}{m+m'}\frac1{m'}\Gamma_{\tau'}^{\mathcal F,m'}
 \leq
 \frac1{m+m'}
 \Gamma_{\frac{m\tau+m'\tau'}{m+m'}}^{\mathcal F,m+m'}.
\]
Fekete's lemma and upper regularization therefore make
\(\tau\mapsto g_\tau^{\mathcal F}\) concave.  Linear boundedness gives
\(\Gamma_\tau^{\mathcal F}=V_\theta\) for \(\tau\ll0\) and
\(-\infty\) above the asymptotic maximal weight.  The Bergman obstacles
are \(\mathcal I\)-good, so their \(P\)- and \(\mathcal I\)-envelopes
agree.  This proves boundedness and the \(\mathcal I\)-model property.

For a smooth test configuration, choose in degree \(m\) an orthonormal basis
adapted to the filtration, with weights \(\lambda_{m,j}\).  The weighted
Bergman ray
\[
 \frac1m\log\sum_j\rme^{t\lambda_{m,j}}
                         |s_{m,j}|_{h^m}^2
\]
differs from the finite Legendre maximum of the degree-\(m\) obstacles by
at most \(m^{-1}\log h^0(X,mL)\), which tends to zero.  Taking usc tail
envelopes identifies the limiting Bergman ray with the Legendre transform
of \(\Gamma^{\mathcal F}\).  Theorem~9.2 of \cite{RWN14} identifies this
limit with the Phong--Sturm ray; uniqueness of the Perron solution gives
the same identification on every compact time interval.
\end{proof}

\begin{example}[A product one-parameter subgroup]
\label{ex:rays-p1-action}
Let \(X=\mathbf P^1\), \(L=\mathcal O(1)\), and let
\(\mathbf G_m\) act with weights \(0,1\) on homogeneous coordinates.
The degree-\(m\) weights are \(0,\ldots,m\).  The resulting filtration
gives a bounded test curve with terminal parameter \(1\); its transform
is the product geodesic ray.  Translation of all weights by \(mc\)
changes the ray by \(ct\), exactly as in
\eqref{eq:rays-test-translation}.
With \(t=-\log|z|^2\), the degree-\(m\) weighted Bergman expression uses
\(\rme^{tj}|s_j|_{h^m}^2\).  Replacing \(t\) by \(-\log|z|\) without
simultaneously changing the weights would therefore produce the wrong
product ray.  This checks the squared-norm convention on the simplest
nonconstant example.
\end{example}

\begin{example}[The hyperplane test curve on \(\mathbf P^n\)]
\label{ex:rays-hyperplane-test-curve}
In the simplest divisorial case every object of this chapter is explicit,
and one computation exercises the test-curve axioms, the Legendre
transform, the linear-supremum lemma and the energy formula at once.

Let \(X=\mathbf P^n\), let \(L=\mathcal O(1)\) carry the Fubini--Study
metric \(h\), and put \(\theta=c_1(L,h)=\omega_{\mathrm{FS}}\), normalized
so that \(\int_X\theta^n=1\).  Let \(H=\{x_0=0\}\) and
\[
 \psi:=\log|s_H|_h^2=\log\frac{|x_0|^2}{|x|^2},
\]
so that \(\theta+\ddc\psi=[H]\) and \(\sup_X\psi=0\), the normalization
demanded before \eqref{eq:rays-linear-supremum}.  Let \(\mathcal F\) filter
\(R(X,L)\) by order of vanishing along \(H\), so that
\[
 \mathcal F^{m\tau}H^0(X,mL)
 =x_0^{\lceil m\tau\rceil}\,
  H^0\bigl(X,\mathcal O(m-\lceil m\tau\rceil)\bigr).
\]

This vanishes exactly when \(\tau>1\), so
\(\Gamma^{\mathcal F}_{\max}=1\).  Write \(a_i=|x_i|^2/|x|^2\), so that
\(a_0=\rme^{\psi}\) and \(\sum_ia_i=1\).  The sup-norm normalization
\(|s|^2_{h^m}\leq1\) in \eqref{eq:rays-bergman-test} is what makes the
computation nontrivial: a monomial \(x^J\) with \(J_0\geq m\tau\) has
\(\sup_X|x^J|^2_{h^m}=\prod_i(J_i/m)^{J_i}\), so the normalized monomial
contributes
\[
 \sum_i\sigma_i\log a_i+H(\sigma),
 \qquad\sigma=J/m\in\Delta_n,\quad
 H(\sigma)=-\sum_i\sigma_i\log\sigma_i ,
\]
an entropy term that a naive count omits.  Maximizing over
\(\sigma\in\Delta_n\) subject to \(\sigma_0\geq\tau\), the unconstrained
maximum being \(0\) at \(\sigma=a\),
\eqref{eq:rays-bergman-test}--\eqref{eq:rays-filtration-test-curve} give
\[
 \Gamma_\tau^{\mathcal F}
 =\begin{cases}
   V_\theta=0, & \tau\leq0\ \text{ or }\ a_0\geq\tau,\\[2pt]
   \tau\log a_0+(1-\tau)\log(1-a_0)+H_2(\tau),
     & a_0\leq\tau\leq1,
  \end{cases}
\]
with \(H_2(\tau)=-\tau\log\tau-(1-\tau)\log(1-\tau)\).  This is
\emph{not} \(\tau\psi\), and the discrepancy is visible at \(m=2\): for
\(n=1\), \(\tau=\tfrac12\), the section \(s=2x_0x_1\) lies in
\(\mathcal F^{1}H^0(\cO(2))\) and has \(|s|^2_{h^2}=4a_0(1-a_0)\leq1\),
so it is admissible, and at \(a_0=\tfrac14\) it gives
\(\tfrac12\log|s|^2_{h^2}=\tfrac12\log\tfrac34=-0.1438\), whereas
\(\tfrac12\psi=\tfrac12\log\tfrac14=-0.6931\).  The displayed formula
returns exactly \(-0.1438\) there.  Each level is still a model potential
with analytic singularities, of coefficient \(\tau\) along \(H\) and nowhere
else: \(\Gamma_\tau-\tau\psi\) equals \((1-\tau)\log(1-a_0)+H_2(\tau)\) on
\(\{a_0<\tau\}\), where \(1-a_0>1-\tau>0\), and \(-\tau\log a_0\) on
\(\{a_0>\tau\}\), where \(a_0>\tau>0\), so it is bounded.  In particular
there is no singularity at \(P_0=[1:0:\cdots:0]\), which lies in the region
\(\{a_0>\tau\}\) on which the level vanishes identically.  So its \(P\)- and
\(\mathcal I\)-envelopes agree and the curve is \(\mathcal I\)-model, as
\Cref{thm:rays-filtration-interface} asserts in general.

The mass is unchanged by the correction.  On \(\{a_0>\tau\}\) one has
\(\theta+\ddc\Gamma^{\mathcal F}_\tau=\theta\); on \(\{a_0<\tau\}\) it is
\(\tau[H]+(1-\tau)\eta\), where \(\eta\) is the pullback of
\(\omega_{\mathrm{FS}}\) under the projection from \(P_0\), so
\(\eta^n=0\); and the nonpluripolar product puts no mass on the pluripolar
set \(H\).  Since \(\theta^n\) pushes forward under \(x\mapsto a\) to the
normalized Lebesgue measure of \(\Delta_n\),
\[
 M_\theta\bigl(\Gamma_\tau^{\mathcal F}\bigr)
 =\int_{\{a_0>\tau\}}\theta^n
 =(1-\tau)^n
 \qquad(0\leq\tau\leq1).
\]
This is positive exactly for \(\tau<1\), which both confirms
\Cref{lem:rays-test-positive-levels} and shows that the terminal parameter
could not have been taken any larger.

The Legendre transform is the associated log partition function, the
supremum over \(\tau\) being attained at
\(\tau=\rme^ta_0/(\rme^ta_0+1-a_0)\geq a_0\):
\[
 \Gamma_t^{\mathcal F*}
 =\sup_{\tau\leq1}\bigl(\Gamma^{\mathcal F}_\tau+t\tau\bigr)
 =\log\bigl(\rme^ta_0+1-a_0\bigr)
 =\log\bigl(1+(\rme^t-1)\rme^{\psi}\bigr),
 \qquad t>0 .
\]
This is the product geodesic ray attached to the deformation of
\(\mathbf P^n\) to the normal cone of \(H\).  It is \emph{not}
\(\max\{0,t+\psi\}\), which is not a geodesic ray at all: for \(n=1\),
where \(\theta\) pushes forward to \(\rd a_0\) on \([0,1]\), one computes
\(E(\max\{0,t+\psi\})=\tfrac12(t-1+\rme^{-t})\), so that
\(E\) at \(t=2\) is \(\tfrac12(1+\rme^{-2})=0.5677\) while twice its value
at \(t=1\) is \(\rme^{-1}=0.3679\) --- not affine, as a weak geodesic must
be --- whereas the displayed ray has \(E(\Gamma^{\mathcal F*}_t)=t/2\)
exactly.  Its supremum is \(t\), attained at \(P_0\) where \(a_0=1\), so
\eqref{eq:rays-linear-supremum} holds with
\(\sup_X\Gamma_1^{\mathcal F*}=1\).

Finally \(\phi=\Gamma_{-\infty}^{\mathcal F}=0\) has
\(M_\theta(\phi)=1\), and \eqref{eq:rays-test-energy} evaluates to
\[
 \mathbf E^\phi(\Gamma^{\mathcal F})
 =1\cdot1+\int_{-\infty}^{1}
    \bigl(M_\theta(\Gamma_\tau^{\mathcal F})-1\bigr)\,\rd\tau
 =1+\int_0^1\bigl((1-\tau)^n-1\bigr)\,\rd\tau
 =\frac1{n+1}.
\]
The integrand vanishes identically for \(\tau\leq0\), which is why the
integral over a half-line converges; this is the mechanism behind the
remark following \eqref{eq:rays-test-energy}.  For \(n=1\) the curve is the
one of \Cref{ex:rays-p1-action}, with its weights and its energy
\(\tfrac12\) now written out.
\end{example}

\begin{remark}[Source note]
\label{rem:rays-source-note}
The analytic correspondence is
\cite[Theorems~1.1--1.2]{RWN14}; the filtration ray and its
Phong--Sturm identification are \cite[Theorem~9.2]{RWN14}.  The relative
and \(\mathcal I\)-model form
used here follows the supplied chapters on rays and test curves and the
transcendental framework of \cite{DXZ23}.  In particular, all test-curve
levels in the notation \(\PSH^{\NA}(X,\theta)_{>0}\) have positive
nonpluripolar mass; the zero-volume objects introduced next are inverse
systems of such perturbed positive-mass components.
\end{remark}
